% Mechanical relocation generated from the preservation ledger.
% Existing prose is included byte-for-byte from preserved blocks.

\section{Conditional Gravity Demand}
\subsection{Model Card: Reduced Conditional Gravity}

\paragraph{Formulation and dimension.}
Within each origin--departure-time production group, the model distributes a
declared journey total across feasible destinations through a conditional
logit expression based on supplied destination attractiveness
(Section~\ref{sec:destination-attractiveness}), scheduled travel time, transfers,
and any explicitly introduced extensions.  Its dimension is the number of
shared behavioral coefficients plus the dimension of a declared production
basis, rather than the number of OD cells.

\paragraph{Required information and identifying assumptions.}
The minimal model requires feasible journey alternatives, their fixed routing
shares and attributes, a positive destination-attractiveness baseline and
correction (Section~\ref{sec:destination-attractiveness}), and either
externally justified journey productions or a low-dimensional production
basis.  It assumes a common parametric response to the included attributes and
conditional independence from omitted attributes.  The normalization within
each production group identifies only relative destination utility; an
unconstrained common additive utility is not identified.

\paragraph{Recoverable quantities and limitations.}
The shared parameters can be identified from count patterns even when
individual OD cells cannot, and the fitted OD table follows from the estimated
parameters and declared productions.  However, this apparent precision is
conditional on the gravity form, destination-attractiveness baseline and
correction, feasible-set
construction, fixed routing, and production assumptions.  Omitted spatial or
temporal heterogeneity can create systematic residuals.  Because its parameter
dimension remains small, this is the preferred primary demand model for the
full-network MAP estimator, subject to adequacy and sensitivity checks.


\subsection{Destination Attractiveness: Baseline, Correction, and Effective Value}
\label{sec:destination-attractiveness}

Destination attractiveness is a supplied input to the public gravity package,
not a quantity that the package infers directly from observed counts.  Let
$g=(o,t)$ denote an origin--departure-time production group and let
$c=(o,d,t)$ denote one OD--time cell.  The set $C_g$ contains the feasible free
cells in group $g$.  The feature provider supplies a strictly positive baseline
destination attractiveness $A^{(0)}_{dt}>0$.  An optional estimated
destination correction is denoted by $\eta_d$; it is shared by all departure
times for destination $d$.  If a destination--time correction is selected, the
correction is instead written $\eta_{dt}$.

For the destination-only case, the effective attractiveness is
\begin{equation}
    A^{\mathrm{eff}}_{dt}
    = A^{(0)}_{dt}\exp(\eta_d).
\end{equation}
With $T_{odt}$ denoting expected scheduled travel time, $K_{odt}$ expected
transfers, and $I_{odt}$ any explicitly declared interaction, the conditional
utility is
\begin{equation}
    V_{odt}
    = \log A^{(0)}_{dt} + \eta_d
      -\beta_T\frac{T_{odt}}{T_{\mathrm{scale}}}
      -\beta_K K_{odt} + I_{odt}.
\end{equation}
The corresponding within-origin--time probability and OD demand are
\begin{equation}
    p_{odt}
    = \frac{\exp(V_{odt})}
      {\displaystyle\sum_{(o,d',t)\in C_{(o,t)}}\exp(V_{od't})},
    \qquad
    x_{odt}=Q_{ot}p_{odt}.
\end{equation}
Here $Q_{ot}$ controls the total demand produced by origin $o$ in time bin
$t$.  Attractiveness and journey attributes control how that total is divided
among feasible destinations; destination attractiveness does not create
additional production.  A destination correction is therefore not
interchangeable with a production correction.

\paragraph{Construction of the baseline.}
The public package receives the positive cell-level vector through
\texttt{GravityFeatures.destination\_attractiveness}.  The feature provider or
case adapter constructs this vector and records its provenance; there is no
universal package-side calculation from observations.  In the current
medium-network adapter, the baseline is constructed as
\begin{equation}
    A^{(0)}_{dt}
      = \sum_{o:(o,d,t)\in C_{\mathrm{free}}}q^{(0)}_{odt},
\end{equation}
where $q^{(0)}_{odt}$ is the baseline flow for a free OD--time cell.  The sum
is over free cells only.  Frozen cells contribute through the separate frozen-
demand offset and are not folded into this baseline.  The resulting value is
copied to every free cell with destination $d$ and time bin $t$.  With the
neutral all-ones seed, this construction is simply the number of eligible free
origins for $(d,t)$.  It is a model input or prior construction, not an
observed destination-demand total, and should not be called ``measured
attractiveness'' unless an external data source actually supplies it.

\paragraph{Fixed and estimated modes.}
The baseline remains supplied in every mode; what changes is whether an
additional correction is estimated:
\begin{table}[htbp]
\centering\small
\setlength{\tabcolsep}{4pt}
\begin{tabular}{p{0.18\linewidth}p{0.22\linewidth}p{0.24\linewidth}p{0.27\linewidth}}
\hline
Mode & Baseline offset & Estimated correction & Interpretation\\
\hline
Fixed & supplied $A^{(0)}_{dt}$ & none & Destination allocation is determined by the baseline and behavioral attributes\\
Estimated & supplied $A^{(0)}_{dt}$ & $\eta_d$ & Counts may adjust relative destination attractiveness\\
Destination--time estimated & supplied $A^{(0)}_{dt}$ & $\eta_{dt}$ & Allows time-varying destination corrections\\
\hline
\end{tabular}
\caption{Destination-attractiveness modes.}
\end{table}
Thus \texttt{mode = "fixed"} does not mean that all destinations have
attractiveness one.  It means that no additional destination correction is
estimated.  The baseline, correction, and their product have distinct roles:
$A^{(0)}_{dt}$ is a supplied multiplicative offset, $\eta_d$ is an estimated
additive utility correction, and $A^{\mathrm{eff}}_{dt}$ is the resulting
effective attractiveness.

\paragraph{Normalization and regularization.}
For a destination correction with a sum-to-zero constraint,
\begin{equation}
    \sum_{d=1}^{D}\eta_d=0.
\end{equation}
Only $D-1$ corrections are stored; the last is reconstructed as the negative
sum of the others.  Without this constraint, adding a common constant to all
$\eta_d$ would leave every conditional probability unchanged because the
constant cancels in the softmax denominator.  The implemented ridge penalty is
\begin{equation}
    R_A(\boldsymbol\eta)
      = \frac{\lambda_A}{2}\sum_{d=1}^{D}\eta_d^2.
\end{equation}
The current case adapter specifies $\lambda_A$ directly through the
\texttt{destination\_ridge} setting in \texttt{GravityRegularization}.  Its
\texttt{prior\_scale} setting is not consumed by this adapter and must not be
described as active.  If a future interface instead specifies a Gaussian prior
standard deviation $\sigma_A$, the equivalent ridge strength is
$\lambda_A=1/\sigma_A^2$, unless that model defines the prior differently.

\paragraph{A small numerical example.}
Suppose $A^{(0)}_{d_1,t}=4$, $A^{(0)}_{d_2,t}=2$,
$\eta_{d_1}=0.1$, and $\eta_{d_2}=-0.1$.  Then
\begin{equation}
    A^{\mathrm{eff}}_{d_1,t}=4e^{0.1},
    \qquad
    A^{\mathrm{eff}}_{d_2,t}=2e^{-0.1}.
\end{equation}
These effective values are combined with the travel-time and transfer
penalties (and any declared interaction) before the within-origin--time
normalization.

\paragraph{Public and case-specific responsibilities.}
The public gravity package defines the mathematical contract and consumes the
feature vector.  The case adapter constructs the baseline vector, the
free/frozen mask, destination and time indices, and provenance from
case-specific OD inputs.  Different cases may therefore use different
positive baseline constructions, provided that they are explicitly documented
and included in model provenance.

\subsection{Minimal Conditional Model}

For feasible cell $c=(o,d,t)$ in production group $g=(o,t)$, let
$d(c)=d$ and $t(c)=t$.  Write $a_c=A^{\mathrm{eff}}_{d(c),t(c)}>0$ for the
effective attractiveness defined above, and let $T_c$ and $K_c$ be expected
scheduled travel time and transfers.  With no estimated destination correction
($\eta_d=0$), the minimal utility, conditional share, and demand are
\[
 V_c=\log a_c-\beta_T T_c/T_{\mathrm{scale}}-\beta_KK_c,
 \qquad
 p_c=\frac{\exp(V_c)}{\sum_{k\in C_g}\exp(V_k)},
 \qquad
 x_c=Q_{ot}p_c.
\]
Externally supplied $Q_{ot}$ gives the smallest model. If productions are
estimated through a declared basis $Z$, a low-dimensional alternative is
\[
 Q_{ot}=Q_{ot}^{(0)}\exp((Z\bm\gamma)_{ot}).
\]
Structural-zero and frozen cells do not receive free gravity-response rows.
The count mean is obtained by applying the response operator to $\bm x$ and
adding the frozen offset.

\subsection{Time-Regime Production Deviations}

The current public implementation provides a specific production extension
when the supplied origin--time totals are systematically too high or too low
in broad, case-defined time regimes.  Let $r(t)$ map each detailed
departure-time bin to one of $R$ regimes.  The production multiplier is
\[
    \log m_{o,t}=\alpha+\delta_{r(t)},
    \qquad
    Q_{ot}=Q^{(0)}_{ot}\,m_{o,t},
    \qquad
    x_{odt}=Q_{ot}p_{odt}.
\]
Here $\alpha$ is the existing global log production scale and $\delta_r$ is a
regime deviation.  The mapping is supplied by the case adapter through the
feature array \texttt{time\_period\_index}; the package does not assume names
such as morning, evening, or peak and does not infer regimes from counts.

For the default sum-to-zero constraint, only $R-1$ deviations are stored and
the final one is the negative sum of the others:
\[
    \sum_{r=1}^{R}\delta_r=0.
\]
This keeps $\alpha$ identifiable as the overall production level.  The
deviation block is ridge-regularized while the global scale is not:
\[
    R_{\mathrm{time}}(\bm\delta)
      =\frac{\lambda}{2}\sum_{r=1}^{R}\delta_r^2,
    \qquad
    J=-\ell_{\mathrm{counts}}-\ell_{\mathrm{auxiliary}}
      +R_{\mathrm{time}}(\bm\delta)+\cdots .
\]
The unit weights shown here are the implemented default.  The penalty is in
the native objective units and is not rescaled by the optimizer's typical
objective scale.  Its gradient is therefore included in the reported raw and
Dennis--Schnabel scaled-gradient diagnostics.

The parameter layout contains the stable names
\texttt{production\_scale}, followed by
\texttt{production\_time\_deviation[0]}, \ldots .  Setting all deviations to
zero reproduces the global-production model exactly; nonzero deviations alter
only the origin--time totals assigned to their mapped regimes, not the
within-group destination normalization.  A change to this gravity
specification changes the gravity-model fingerprint and requires a new gravity
fit checkpoint, but it does not change the routing or fixed-operator artifact
identity, so those expensive artifacts remain reusable.

\subsection{Declarative Gravity Specifications}

The implementation no longer assumes that the first three parameters exhaust
the model. A versioned \texttt{GravityModelSpecification} declares the scope,
parameterization, normalization, feature mapping, and regularization of every
component. The conservative template is
\path{configs/gravity/gravity_model_spec.yaml}. It can be loaded and validated
without editing estimation code; unknown fields and unsupported combinations
are errors rather than ignored options.

The supported categorical scopes are global, origin, destination, time period,
origin--time, destination--time, origin zone, destination zone, zone pair, and
a named custom group. The scopes \texttt{none} and \texttt{fixed} introduce no
estimated parameter. A smooth temporal basis is also supported when its
cell-by-basis matrix is present in the feature cache and is named in the time
specification. These scopes can be applied, subject to component-specific
validation, to production corrections, destination-attractiveness corrections,
journey-time, transfer, and waiting-time coefficients, and departure-time
effects. Negative-binomial dispersion is global or fixed. Estimated residual
demand is deliberately rejected because no supported observation model yet
identifies it separately.

For a general declared model, cell utility and production are
\[
\begin{aligned}
 V_c & = \log a_c + \alpha_c + \tau_c
       -\beta_{T,c}T_c/T_{\mathrm{scale}}
       -\beta_{K,c}K_c-\beta_{W,c}W_c,\\
 q_g & =q_g^{(0)}\exp(\delta_g),
 & x_c & =q_{g(c)}p_c.
\end{aligned}
\]
In this abstract form, $a_c$ is the effective attractiveness
$A^{\mathrm{eff}}_{d(c),t(c)}$ supplied by the baseline and any selected
correction.  More explicitly, the destination contribution is
$\log A^{(0)}_{d(c),t(c)}+\eta_{d(c)}$ for a destination correction, or
$\log A^{(0)}_{d(c),t(c)}+\eta_{d(c),t(c)}$ for a destination--time correction;
$\tau_c$ denotes a separate temporal utility contribution.  The feature cache
provides $A^{(0)}$ through
\texttt{features.destination\_attractiveness}, while
\texttt{destination\_index} and \texttt{time\_period\_index} provide the
grouping used by the declared
\texttt{GravityComponent\allowbreak{}Specification} and
\texttt{GravityParameter\allowbreak{}Layout.cell\_effect}.  Thus production corrections
change $q_g$, whereas destination corrections change relative shares within
$g$.

Positive behavioral coefficients use a softplus transformation with a small
strictly positive floor. Production corrections are log multipliers;
attractiveness and temporal corrections are additive utilities. Categorical
deviations use either an explicit reference category or an exact sum-to-zero
expansion. A grouped positive coefficient consists of one positive base and
normalized log deviations. The implemented time-regime production deviation
is a separate centered block attached to a global production scale and must
carry ridge regularization. More general origin--time and custom production
corrections must also carry ridge regularization because their dimension can
otherwise approach the number of production groups.

The parameter layout is derived deterministically from the specification. It
records one named slice per block, raw-to-physical and physical-to-raw
transformations, constraints, regularization, and a stable fingerprint.
Legacy names \texttt{beta\_time}, \texttt{beta\_transfer},
\texttt{dispersion}, and \texttt{production\_scale} are retained, so the
original three-parameter model and the optional global production scale remain
backward compatible. Nested specifications are warm-started by matching these
stable parameter names; new normalized deviations start at zero and reproduce
the parent prediction exactly.

\subsection{Feature and Identifiability Contracts}

\texttt{GravityFeatures} is an immutable sparse cell table, not a dense
origin--destination--time tensor. In addition to the canonical origin,
destination, departure-time, and origin--time indices, it may contain time
period, origin-zone, destination-zone, destination--time, zone-pair, named
custom-group, waiting-time, and smooth-time-basis arrays. Required mappings
are inferred from the specification and checked before estimation for length,
integer type, contiguous labels, declared group count, and constancy within a
production group where required. All optional mappings participate in the
feature-cache fingerprint.

Invalid specifications are rejected, including a global production scale with
an estimated detection-rate scale, an unnormalized categorical effect, and an
unregularized high-dimensional production correction. Technically valid but
weakly identified choices produce visible warnings. Examples are a global
additive attractiveness correction, which cancels in each conditional softmax,
and simultaneous grouped production and attraction corrections, which may be
strongly correlated even after normalization. Model selection must therefore
use grouped holdout performance, gradient and curvature diagnostics,
regularization sensitivity, and count-mass audits rather than calibration
likelihood alone.

The generalized generator preserves structural zeros exactly and applies only
sparse per-cell vectors and grouped segment reductions. Its JAX implementation
and independent NumPy reference agree at the established floating-point
tolerances. The same demand function is exercised through negative-binomial
objectives, calibration masks, fixed-positive measurement offsets, batched
forward gradients, and matrix-free adjoint gradients. Unsupported measurement
rows must be excluded by an explicit calibration mask.

\subsection{Information Content and Measurement Design}

The gravity model can diagnose weak information about its own parameters.
Sensitivities of predicted counts, likelihood curvature, profile behavior,
prior sensitivity, and posterior uncertainty reveal parameter combinations
that existing measurements distinguish poorly. Measurement-level derivatives
also show which observed rows contribute to those combinations. This is not
the same as residual analysis: a large residual indicates inadequate fit,
whereas weak curvature indicates inadequate information, and either can occur
without the other.

The low parameter dimension does not imply that every reconstructed OD cell is
locally informed by counts. A precise-looking cell estimate may follow mainly
from the gravity form, attractiveness, production, and feasible-set assumptions.
Choosing new counter locations is therefore a separate experimental-design
problem. Candidate measurements must be compared by their expected increase in
information rank or conditioning, or by the expected reduction in uncertainty
for declared quantities of interest. The gravity response supplies the
derivatives needed for that analysis, but it does not automatically select new
measurement locations.
